% Deutsche Parallelfassung zu geometry/curvature_classes.tex

\chapter{Krümmungsklassen}
\label{chap:curvature-classes}

\section*{Motivation}

Nicht alle Krümmungsgesetze besitzen dieselben geometrischen, dynamischen,
numerischen oder physikalischen Eigenschaften. Innerhalb von AIM werden sie
daher nach Stetigkeit, Differenzierbarkeit, Beschränktheit,
Übergangsverhalten, Realisierungsverhalten, Laufzeitstabilität und
betrieblicher Eignung klassifiziert.

Diese Klassifikation bildet die Grundlage für Übergangsbewertung,
Laufzeitrobustheit, Realisierungsanalyse, Zulässigkeitsprüfung und
Optimierungsstrategien.

\begin{principle}[Krümmungsklassifikationsprinzip]
\label{princ:curvature-classification}
Innerhalb von AIM werden Eisenbahn-Alignment-Elemente primär nach ihrem
Krümmungsverhalten und nicht nach ihrer sichtbaren Form im Positionsraum
klassifiziert.
\end{principle}

\section{Grundlegende Krümmungsklassen}

\subsection{Konstante Krümmung}

\begin{definition}[Konstante Krümmung]
\label{def:constant-curvature-class}
Ein Krümmungsgesetz ist konstant, wenn
\[
\kappa(s)=\kappa_0.
\]
\end{definition}

Sonderfälle sind \(\kappa_0=0\) für eine Gerade und \(\kappa_0\neq 0\) für
einen Kreisbogen. Konstante Krümmung bildet die Grundlage fester Elemente der
klassischen Eisenbahngeometrie.

\subsection{Stückweise konstante Krümmung}

\begin{definition}[Stückweise konstante Krümmung]
\label{def:piecewise-constant-curvature}
Ein Krümmungsgesetz ist stückweise konstant, wenn auf Teilintervallen \(I_i\)
\[
\kappa(s)=\kappa_i
\]
gilt.
\end{definition}

Dies entspricht der klassischen Tangenten--Bogen-Geometrie der Eisenbahn.
Stückweise konstante Krümmung ist geometrisch einfach, Krümmungssprünge sind
jedoch dynamisch und betrieblich problematisch.

\section{Stetigkeitsklassen}

\subsection{\texorpdfstring{$C^0$}{C0}-Krümmung}
\begin{definition}[\(C^0\)-stetige Krümmung]
\label{def:c0-curvature}
Ein Krümmungsgesetz gehört zu \(C^0\), wenn \(\kappa(s)\) stetig ist.
\end{definition}
Stetige Krümmung beseitigt Krümmungssprünge.

\subsection{\texorpdfstring{$C^1$}{C1}-Krümmung}
\begin{definition}[\(C^1\)-stetige Krümmung]
\label{def:c1-curvature}
Ein Krümmungsgesetz gehört zu \(C^1\), wenn \(\kappa'(s)\) existiert und
stetig ist.
\end{definition}
Dies verbessert dynamische Glätte, Realisierungsverhalten und numerische
Robustheit.

\subsection{\texorpdfstring{$C^2$}{C2}-Krümmung}
\begin{definition}[\(C^2\)-stetige Krümmung]
\label{def:c2-curvature}
Ein Krümmungsgesetz gehört zu \(C^2\), wenn \(\kappa''(s)\) existiert und
stetig ist.
\end{definition}

Höhere Stetigkeit reduziert Anregungseffekte, numerische Instabilität und
unerwünschtes hochfrequentes Verhalten. In AIM ist die Stetigkeitsklasse nicht
nur eine mathematische Eigenschaft, sondern ein Indikator für dynamische
Qualität, Laufzeitstabilität und Realisierungsrobustheit.

\section{Beschränktheitsklassen}

\subsection{Beschränkte Krümmung}
\begin{definition}[Beschränkte Krümmung]
\label{def:bounded-curvature}
Ein Krümmungsgesetz ist beschränkt, wenn
\[
|\kappa(s)|\le\kappa_{\max}.
\]
\end{definition}
Beschränkte Krümmung entspricht einem minimal zulässigen Radius.

\subsection{Beschränkter Krümmungsgradient}
\begin{definition}[Beschränkter Krümmungsgradient]
\label{def:bounded-curvature-gradient}
Ein Krümmungsgesetz besitzt einen beschränkten Krümmungsgradienten, wenn
\[
|\kappa'(s)|\le\eta_{\max}.
\]
\end{definition}
Dies steht in engem Zusammenhang mit Ruckbegrenzungen in der
Eisenbahndynamik.

\subsection{Beschränkte höhere Ableitungen}
\begin{definition}[Beschränkte höhere Ableitungen]
\label{def:bounded-higher-curvature-derivatives}
Beschränktheitsbedingungen höherer Ordnung umfassen
\[
|\kappa''(s)|<\infty,\qquad |\kappa'''(s)|<\infty.
\]
\end{definition}
Sie werden für Realisierungsverhalten, Fahrzeugreaktion und den Vergleich von
Übergangsfamilien relevant.

\section{Übergangsklassen}

\subsection{Monotone Übergänge}
\begin{definition}[Monotoner Übergang]
\label{def:monotone-transition}
Ein Übergangskrümmungsgesetz ist monoton, wenn \(\kappa'(s)\) sein Vorzeichen
nicht wechselt.
\end{definition}
Monotone Übergänge repräsentieren eine direkte Krümmungsänderung zwischen zwei
Krümmungszuständen.

\subsection{Symmetrische Übergänge}
\begin{definition}[Symmetrischer Übergang]
\label{def:symmetric-transition}
Ein normiertes Übergangsgesetz ist symmetrisch, wenn sein Krümmungsverlauf
bezüglich der Mitte des Übergangsintervalls symmetrisch ist.
\end{definition}
Für normierte Übergangskoordinaten \(w\in[0,1]\) lässt sich dies durch eine
Symmetriebeziehung der Übergangsformfunktion ausdrücken. Symmetrie ist keine
absolute Anforderung, aber eine nützliche Klassifikationseigenschaft.

\subsection{Asymmetrische Übergänge}
\begin{definition}[Asymmetrischer Übergang]
\label{def:asymmetric-transition}
Ein Übergangsgesetz ist asymmetrisch, wenn die Symmetrie bewusst verletzt
wird.
\end{definition}
Asymmetrische Übergänge können in beschränkten Ingenieursituationen, bei
realisierungsbewusster oder fahrzeugdynamischer Optimierung und in
topologisch beschränkten Alignments auftreten.

\section{Analytische Klassen}

\subsection{Polynomiale Krümmungsgesetze}
\begin{definition}[Polynomiale Krümmung]
\label{def:polynomial-curvature}
Ein polynomiales Krümmungsgesetz erfüllt
\[
\kappa(s)=\sum_{i=0}^n a_i s^i.
\]
\end{definition}
Polynomiale Gesetze sind nützlich, weil Ableitungen und Integrale häufig
systematisch behandelt werden können.

\subsection{Trigonometrische Krümmungsgesetze}
\begin{definition}[Trigonometrische Krümmung]
\label{def:trigonometric-curvature}
Ein trigonometrisches Krümmungsgesetz enthält sinus- oder kosinusbasierte
Anteile.
\end{definition}
Solche Gesetze sind für glattes Übergangsverhalten und frequenzorientierte
Interpretation nützlich.

\subsection{Lookup-basierte Krümmungsgesetze}
\begin{definition}[Lookup-basierte Krümmung]
\label{def:lookup-based-curvature}
Ein Lookup-basiertes Krümmungsgesetz wird durch abgetastete oder tabellierte
Repräsentationen zusammen mit Interpolationsregeln definiert.
\end{definition}
So lassen sich Übergangsfamilien oder Realisierungseffekte praktisch
einbinden, die nicht bequem durch eine einfache geschlossene Formel
darstellbar sind.

\subsection{Hybride Krümmungsgesetze}
\begin{definition}[Hybrides Krümmungsgesetz]
\label{def:hybrid-curvature-law}
Ein hybrides Krümmungsgesetz kombiniert parametrische Struktur, abgetastete
Korrekturen, Messdaten und lokale Verfeinerung.
\end{definition}
Hybride Gesetze sind besonders für gemessene und realisierte
Eisenbahnzustände wichtig.

\section{Dünn besetzte Krümmungsklassen}

\begin{definition}[Dünn besetzte Krümmungsrepräsentation]
\label{def:sparse-curvature-class}
Ein dünn besetztes Krümmungsmodell repräsentiert Geometrie durch feste
Krümmungselemente, Übergangselemente und kompakte Parametersätze.
\end{definition}

Dünn besetzte Klassen betonen Kompaktheit, Laufzeiteffizienz,
Rekonstruktionsstabilität und semantische Klarheit. Innerhalb von AIM sind sie
nicht bloß Speicherformate, sondern definieren ausführbare
Krümmungsprogramme.

\section{Betriebliche Klassen}

\subsection{Dynamisch glatte Krümmung}
\begin{definition}[Dynamisch glatte Krümmung]
\label{def:dynamically-smooth-curvature}
Ein Krümmungsgesetz ist dynamisch glatt, wenn der Ruck beschränkt bleibt, die
Überhöhungsrate zulässig ist, die Krümmungsänderung kontrolliert wird und
hochfrequente Anregung begrenzt ist.
\end{definition}

\subsection{Realisierungsstabile Krümmung}
\begin{definition}[Realisierungsstabile Krümmung]
\label{def:realization-stable-curvature}
Ein Krümmungsgesetz ist realisierungsstabil, wenn die Realisierung seine
dominante Struktur bewahrt, Bauprozesse keine Instabilität einführen und sein
niederfrequentes Verhalten nach der physischen Realisierung identifizierbar
bleibt.
\end{definition}
Realisierungsstabilität ist eine physikalische, nicht bloß eine analytische
Glattheitseigenschaft.

\subsection{Laufzeitstabile Krümmung}
\begin{definition}[Laufzeitstabile Krümmung]
\label{def:runtime-stable-curvature}
Ein Krümmungsgesetz ist laufzeitstabil, wenn Projektion numerisch stabil,
Frame-Bewertung stetig und Pose-Rekonstruktion robust bleibt und lokale
inverse Probleme gut konditioniert bleiben.
\end{definition}

\section{Frequenzklassen}

\subsection{Niederfrequente Krümmung}
\begin{definition}[Niederfrequente Krümmung]
\label{def:low-frequency-curvature}
Niederfrequente Krümmungsanteile bestimmen das großräumige geometrische
Verhalten.
\end{definition}
Diese Anteile werden üblicherweise durch Bau- und Instandhaltungsprozesse
bewahrt.

\subsection{Hochfrequente Krümmung}
\begin{definition}[Hochfrequente Krümmung]
\label{def:high-frequency-curvature}
Hochfrequente Krümmungsanteile beeinflussen
Realisierungsempfindlichkeit, dynamische Rauheit, Messempfindlichkeit und
numerische Instabilität.
\end{definition}

\begin{proposition}[Spektrales Realisierungsprinzip]
\label{prop:curvature-classes-spectral-realization}
Hochfrequente Krümmungsanteile werden durch Realisierung abgeschwächt.
\end{proposition}
Diese frequenzorientierte Klassifikation verbindet Krümmungsklassen mit der
Interpretation der Realisierung als Filter.

\section{Zulässigkeit}

\begin{definition}[Zulässiges Krümmungsgesetz]
\label{def:admissible-curvature-law}
Ein Krümmungsgesetz ist zulässig, wenn es geometrische,
realisierungsbezogene, betriebliche und dynamische Nebenbedingungen sowie
Bedingungen der numerischen Stabilität erfüllt.
\end{definition}

Zulässigkeit hängt von Fahrzeugdynamik, Realisierungsprozessen, Normen,
Topologie, Laufzeitanforderungen und beabsichtigter betrieblicher Nutzung ab.
Sie ist damit kontextabhängig. Ein Krümmungsgesetz kann analytisch gültig,
aber betrieblich ungeeignet oder unter Realisierung physikalisch instabil sein.

\section{Klassifikation und Qualität}

Klassifikation impliziert nicht automatisch Qualität. Ein sehr glattes
Krümmungsgesetz kann ungeeignet sein, wenn es schwer realisierbar, dynamisch
ungünstig, numerisch instabil, mit der Topologie unvereinbar oder für den
beabsichtigten Betrieb ungeeignet ist.

Innerhalb von AIM wird Krümmungsqualität durch das Zusammenwirken von Klasse,
Kontext, Realisierung und Laufzeitverhalten bewertet.

\section{Kanonische Interpretation}

\begin{proposition}
Verschiedene Krümmungsklassen entsprechen grundlegend verschiedenen
geometrischen, betrieblichen, laufzeitbezogenen und
Realisierungsverhaltensweisen.
\end{proposition}
\begin{proposition}
Innerhalb von AIM wird die Qualität eines Eisenbahn-Alignments primär im
Krümmungsraum und nicht im Positionsraum charakterisiert.
\end{proposition}
\begin{proposition}
Übergangsfamilien werden als strukturierte Teilmengen von Krümmungsklassen
interpretiert.
\end{proposition}
\begin{proposition}
Zulässigkeit ist keine rein geometrische Eigenschaft, sondern eine kombinierte
geometrische, dynamische, laufzeitbezogene und Realisierungseigenschaft.
\end{proposition}

\section{Verbindung zu AIM}

\begin{summary}
Krümmungsklassen bilden das strukturelle Klassifikations-Framework der
AIM-Geometrie. Sie verbinden Stetigkeit, Differenzierbarkeit, Beschränktheit,
Übergangs- und Realisierungsverhalten, Laufzeitstabilität, dynamische Glätte
und Optimierungseignung.

Innerhalb von AIM werden Übergangsbögen daher nicht nur als geometrische
Entitäten, sondern als Mitglieder bestimmter Krümmungsklassen mit
unterschiedlichen betrieblichen Eigenschaften interpretiert. Dies schafft die
Brücke von intrinsischer Geometrie zu Klassifikation, Realisierungsanalyse,
Laufzeitbewertung und Optimierung.
\end{summary}
